\begin{definition}
Let $n\ge1$ and let $\Omega\subseteq\R^n$.  We say that $\Omega$ is a
\emph{bounded Lipschitz domain}, written
$\mathrm{IsBoundedLipschitzDomain}(n,\Omega)$, if
\begin{enumerate}
\item $\Omega$ is open;
\item $\Omega$ is connected and nonempty --- together with (1) this is
  ``$\Omega$ is a domain'';
\item $\Omega$ is bounded; and
\item \emph{(the Lipschitz boundary condition)} for every $x\in\partial\Omega$
  there are $r>0$, $L\ge0$, a unit vector $\nu\in\R^n$ and a function
  $\varphi\colon\R^n\to\R$ such that
  \begin{itemize}
  \item $\varphi$ is $L$-Lipschitz on $\R^n$;
  \item $\varphi(y+t\nu)=\varphi(y)$ for every $y\in\R^n$ and every $t\in\R$;
    and
  \item for every $y\in B_r(x)$,
    \[
    y\in\Omega\quad\Longleftrightarrow\quad \varphi(y-x)<(y-x,\nu),
    \]
    where $(\,\cdot\,,\,\cdot\,)$ is the inner product of $\R^n$, so
    $(y-x,\nu)=\sum_{i=1}^n(y-x)^i\nu^i$.
  \end{itemize}
\end{enumerate}

\noindent\textbf{What clause (4) says, and why it is stated this way.}  Fix
$x\in\partial\Omega$ and the data $r,L,\nu,\varphi$ it provides.  The second
bullet says that $\varphi$ does not change along the direction $\nu$, so
$\varphi$ is a function of the component of its argument orthogonal to $\nu$
alone: writing $y-x=w+s\nu$ with $(w,\nu)=0$ and $s=(y-x,\nu)$, one has
$\varphi(y-x)=\varphi(w)$, and the third bullet reads
$y\in\Omega\iff s>\varphi(w)$.  That is exactly the usual condition that, in the
orthonormal coordinates adapted to $\nu$ at $x$, the set $\Omega$ coincides
inside $B_r(x)$ with the region strictly above the graph of a Lipschitz function
of the remaining $n-1$ coordinates.  Stating $\varphi$ on all of $\R^n$ together
with the invariance $\varphi(y+t\nu)=\varphi(y)$ is equivalent to stating it on
the hyperplane $\nu^{\perp}$, and it avoids introducing coordinates on that
hyperplane; the orientation is fixed by taking $\nu$ to point into $\Omega$,
which costs nothing since $\nu$ is existentially bound and $-\nu$ is available.

\noindent\textbf{Where the paper asks for this.}  Section 4 fixes the convention
at l.~1130, ``Fix now $\Omega'$ a Lipschitz domain such that
$\overline{\Omega'}\subset\Omega$'', and Theorems~Theorem~3 of the thesis, Theorem~5 of the thesis and
Theorem~6 of the thesis each restate it --- l.~1274, l.~1431 and l.~1551, ``Let $\Omega'$ be
a Lipschitz domain such that $\overline{\Omega'}\subset\Omega$''.  The cited
result (C5) carries it too, at l.~264: ``if $\Omega'$ is a bounded Lipschitz
domain, $n<\lambda<n+p$ and $u\in\mathcal L^{p,\lambda}(\Omega')$, then
\ldots''.  Two things follow, and they pull in opposite directions, which is why
this predicate is needed on both sides.  A statement of (C5) \emph{without} the
hypothesis is stronger than the literature, which the Note forbids at l.~258
(``stated as in the literature and not in a stronger form''); Campanato's
theorem is false on domains with an outward cusp, where the measure of
$\Omega'\cap B_\rho(x)$ can decay faster than $\rho^n$.  And a statement of
Theorem~Theorem~3 of the thesis, Theorem~5 of the thesis or Theorem~6 of the thesis without it claims the
conclusion for \emph{arbitrary} bounded open $\Omega'$, which is stronger than
the theorem the paper states; that is a target-faithfulness defect of its own,
independent of (C5).

\noindent\textbf{Boundedness is listed separately although (3) is implied by
nothing else here.}  Clause (3) is the ``bounded'' of (C5)'s own wording; the
paper's \S4 convention gets it instead from $\overline{\Omega'}\subset\Omega$
with $\Omega'$ a domain only when $\Omega$ is itself bounded, which the paper
does not assume, so it is stated.  Clauses (1) and (3) duplicate hypotheses that
the consumers of this predicate already carry by name --- they were introduced
before this predicate existed, and the many proof sites that discharge
``$\Omega'$ is open'' and ``$\Omega'$ is bounded'' individually are left as they
are; the duplication is harmless, being a repeated conjunct, and avoids
rewriting those sites.
\end{definition}
